\appendix
\section{Analytic dependencies}\label{sec:sources}

The results are conditional on the following inputs, with the hypotheses
and normalizations stated in Section~\ref{sec:framework}. References
\cite{D,S,P,Z} refer throughout to the same fixed repository revision.\footnote{Revision
\texttt{fd4aeeb2ee4fc729c18d98444fed42fd0529eeeb} of
\href{https://github.com/openai/math/tree/fd4aeeb2ee4fc729c18d98444fed42fd0529eeeb}{\texttt{openai/math}}.
The bibliography links to these versions.}
Proof estimates mentioned below are included among the assumptions.

\subsection*{Transference and marked correlations}
We use the general graph-transference and residual ideal-family
propositions in \cite[Section 6]{D}, which apply the local transference,
endpoint pairing, and ideal-kernel results of \cite[Sections 3--4]{S}.
The required endpoint norms, list normalization, shared-label collision
bound, and bilinear multiplier estimate are those used in their proofs.
The local Mellin-energy lemma, truncated divisor-model lemma, and sparse
Gram bound in \cite[Section 6]{D} give the reduction
\eqref{eq:energy} and expansion \eqref{eq:divisor}.
In particular, the norm exponent is independent of the fixed shared-list
parameter $J$. Sections~\ref{sec:parity} and~\ref{sec:model} verify the
new endpoint energy estimates; no signed or fixed-progression version
of the final distribution theorem is assumed.

\subsection*{Oscillation and the scalar estimate}
The long-prime polynomial and logarithmic-phase lemmas in
\cite[Appendix A]{D} are used in the ranges
\eqref{eq:longprime}--\eqref{eq:logphase}, including nonunit residue
classes in the latter. For the determinant replacement, we use the
coefficient criterion for the marked estimate in \cite[Section 8]{P}
and the scalar Cauchy-square bound in its proof. Its coefficient
hypothesis is \eqref{eq:coefficient}; the prime-minus-rough coefficient
satisfies it by the minor-moment and major-term estimates in
\cite[Sections 3--5]{D}. Equation~\eqref{eq:detphase} keeps the new
multiplicative factor outside that scalar square.

\subsection*{Positive extraction}
We use the block-sieve and divisor-moment estimates in
\cite[Section 2]{D}, followed by the progression, prime-to-rough,
conditional-coin, allocation, and mean-square deweighting estimates in
\cite[Section 7]{D}. These include the proof estimates for removing
candidate prime parts from the rough coefficients and thresholds.
The sieve level, positive proxy mass, and deweighting bound used here
are stated in Sections~\ref{sec:framework} and~\ref{sec:extraction}.
The complex progression remainder and the recombination of its
principal terms are proved in Section~\ref{sec:extraction}.

\subsection*{Ordinary multiplicative-function estimates}
The main theorem of \cite{Z} supplies the zero-free half-plane for
Dirichlet $L$-functions required in Section~\ref{sec:model}. We use
this assertion as a hypothesis. The remaining distribution inputs are
the Liouville estimate recorded in \cite[equation (1.5)]{Pintz},
Granville--Shao's criterion \cite[Theorem 1.1(b)]{GS}, and the
Selberg--Delange expansion \cite[Theorem 1.2 and equation (1.10)]{BT},
as applied in Sections~\ref{sec:parity} and~\ref{sec:model}.
